\documentclass[a4paper,12pt]{ctexart}

% ========== Packages ==========
\usepackage{amsmath, amssymb, amsthm}
\usepackage{graphicx}
\usepackage{booktabs}
\usepackage{geometry}
\usepackage{hyperref}
\usepackage{listings}
\usepackage{xcolor}
\usepackage{caption}
\usepackage{subcaption}
\usepackage{float}
\usepackage{enumitem}
\usepackage{url}

% ========== Layout ==========
\geometry{margin=2.5cm}
\setlength{\parindent}{2em}
\hypersetup{colorlinks=true, linkcolor=black, urlcolor=blue, citecolor=blue}

% ========== Code listing style ==========
\definecolor{codegray}{rgb}{0.5,0.5,0.5}
\definecolor{codebg}{rgb}{0.97,0.97,0.97}
\lstset{
    backgroundcolor=\color{codebg},
    basicstyle=\ttfamily\footnotesize,
    numbers=left,
    numberstyle=\tiny\color{codegray},
    frame=single,
    breaklines=true,
    showstringspaces=false,
    tabsize=4
}

% ========== Title ==========
\title{\textbf{2026 XUPT Campus Mathematical Contest in Modeling\\ \large Problem X: [Title Here]}}
\author{Team Number: \underline{\quad\quad\quad\quad}}
\date{\today}

\begin{document}

\maketitle

% ========== Abstract ==========
\begin{abstract}
\noindent\textbf{Abstract:} Briefly describe the problem background, methodology, key results, and conclusions in one paragraph (approximately 300--500 words).

\vspace{0.5em}
\noindent\textbf{Keywords:} keyword1; keyword2; keyword3
\end{abstract}

\newpage
\tableofcontents
\newpage

% ==========================================
\section{Problem Restatement}
\subsection{Background}
Restate the problem background concisely.

\subsection{Problem Analysis}
\begin{itemize}
    \item \textbf{Problem 1:} ...
    \item \textbf{Problem 2:} ...
\end{itemize}

% ==========================================
\section{Assumptions and Notations}
\subsection{Fundamental Assumptions}
\begin{enumerate}
    \item Assumption 1.
    \item Assumption 2.
\end{enumerate}

\subsection{Notations}
\begin{table}[H]
    \centering
    \caption{Symbol Descriptions}
    \begin{tabular}{@{}clc@{}}
        \toprule
        Symbol & Description & Unit \\
        \midrule
        $x$ & Decision variable & -- \\
        $t$ & Time & s \\
        $\lambda$ & Lagrange multiplier & -- \\
        \bottomrule
    \end{tabular}
\end{table}

% ==========================================
\section{Model Establishment and Solution}

\subsection{Problem 1: Model Building}
A sample equation:
\begin{equation}
    \min_{x \in \mathbb{R}^n} f(x) \quad \text{s.t.} \quad g_i(x) \le 0,\; i=1,\dots,m
    \label{eq:opt}
\end{equation}

\subsection{Problem 1: Algorithm Design}
Describe the solver or heuristic used.

\subsection{Problem 1: Results}
% \begin{figure}[H]
%     \centering
%     \includegraphics[width=0.7\textwidth]{figures/result1.png}
%     \caption{Result of Problem 1}
%     \label{fig:res1}
% \end{figure}

% ==========================================
\section{Sensitivity Analysis}
Analyze how the solution changes with perturbations to key parameters.

% ==========================================
\section{Model Evaluation}
\subsection{Strengths}
\begin{itemize}
    \item ...
\end{itemize}

\subsection{Weaknesses}
\begin{itemize}
    \item ...
\end{itemize}

\subsection{Extensions}
Potential improvements and generalizations.

% ==========================================
\section{Conclusion}
Summarize findings.

% ==========================================
\begin{thebibliography}{99}
    \bibitem{ref1} Author A. \textit{Title of Paper}. Journal Name, Year.
    \bibitem{ref2} Author B. \textit{Book Title}. Publisher, Year.
\end{thebibliography}

% ==========================================
\appendix
\section{Source Code}
\lstinputlisting[language=Python, caption={Main solver}]{code/main.py}

\end{document}
